\section{The design challenges of autonomous systems}
\label{sec:prb_design_challenges}
\subsection{Defining the tasks}
\label{sec:design_challenges_tasks}
The initial decisions in the design of an autonomous system concern the specifications of the operations that it performs. These specifications are fully defined from the onset down to the smallest detail possible at this stage, as they guide all subsequent choices, from hardware selection and construction to programming. During the breakdown of operational requirements into \emph{tasks}, intended as discrete units of work to be executed, various types emerge. Without delving, for now, into the specifics of a particular scenario, common categories of tasks include the following ones.

\begin{itemize}
    \item \textbf{Continuous tasks}, which begin at system startup and run uninterrupted until it is turned off or malfunctions.
    \item \textbf{Periodic tasks}, which are performed at regular intervals.
    \item \textbf{Asynchronous tasks}, which are executed only when a specific event occurs, which may or may not ever occur.
\end{itemize}

It is easy to see that most functionalities in modern autonomous systems can be broken down into tasks that fall within these categories. These include operations ranging from actuator control, to the collection of data and measurements required by control algorithms, to communication with operators and users. The system architecture should efficiently support the development of software modules and the integration of hardware solutions that can handle all these kinds of tasks.

\subsection{Choosing the hardware}
\label{sec:design_challenges_hardware}
The choices made initially in this direction may not be definitive, because they concern both the construction of the initial prototypes and the final product: there likely are misjudgments or unforeseen issues to address. Nevertheless, the hardware should be selected carefully based on the operational, structural, and construction requirements defined so far. Decisions taken in the previous step are of critical importance at this stage: establishing the computational complexity of the operations that the system has to perform dictates the kind of resources required from the hardware, which may need to be equipped with discrete coprocessors to handle specific types of data or subdivided into multiple interconnected subsystems. The term \emph{discrete} here refers to auxiliary hardware installed in a computer that cannot operate independently but is designed to be driven by the primary computer, providing specialized support for specific tasks.

A market survey is essential to assess whether pre-tested and ready-to-use components or SoCs (System-on-Chip) can be used rather than resorting to a proprietary solution, which generally entails significant development, testing, and certification costs. Another domain that adds its own level of complexity is that of lower-level hardware, which consists of:
\begin{itemize}
    \item high-speed \emph{microcontrollers};
    \item actuators;
    \item sensors;
    \item interfaces and communication devices;
    \item data storage systems.
\end{itemize}

Nothing more can be said about these elements without reference to a specific case. However, practical scenarios discussed later in this work follow this design practice. From the multiplicity of devices mentioned and that may compose a modern autonomous system, it starts to become clear how the complexity of the system can grow quickly, as anticipated in \Cref{ch:intro}. The need for a well-structured architecture becomes more evident.

\subsection{Organizing the software}
\label{sec:design_challenges_software}
Except for those components that perform crucial and highly specific functions, such as electrical or electronic circuits, nearly all operations carried out by a modern autonomous system are encoded within an internal computer, which translates these instructions into commands for other components. Given the underlying complexity of the system itself, the software that controls it also has a layered and highly hierarchical organization. In this hierarchy, the lower levels are occupied by modules that interact with the smaller and faster subsystems (\emph{e.g.}, the \emph{firmware}) and these modules are developed with particular attention to their characteristics and purposes. Moving up in the hierarchy, more modules are found that encode sophisticated operations such as supervision, planning, and user communication, making them closer to the notion of \emph{application} or \emph{system software}.

When operating in this context, it is essential to rely on tools that simplify design and development as much as possible at all levels, assisting with the management of issues related to module organization, inter- and intra-process communication, data processing, and storage of various data types. The convenience of a good architecture to handle these challenges becomes increasingly evident.

\subsection{Handling data}
\label{sec:design_challenges_data}
An autonomous system requires a nearly constant stream of information to complete its tasks. This information can be received in data packets managed by suitable communication interfaces, or it can be acquired directly through measurement devices commonly known as \emph{sensors}. For both solutions, typical challenges that arise immediately include the integration of these devices with the rest of the apparatus and the processing of the data they produce. These issues are crucial, as the operation of the entire system is almost always contingent upon obtaining information about measurable quantities of interest and the system's own state (consider, \emph{e.g.}, a feedback control algorithm).

It is therefore reasonable to expect that the hardware architecture of a modern autonomous system comprises numerous communication interfaces and sensors, the latter consisting of specific hardware designed to measure particular physical quantities of interest. The challenges mentioned above concern the integration of such hardware with the rest of the system and to the development of software modules capable of managing it, and of receiving and processing the data it generates.

\subsection{Debugging and testing}
\label{sec:design_challenges_testing}
As soon as the initial prototypes are finalized, testing should begin to verify functionality and confirm the validity of the design decisions that have been taken. Bearing in mind that a generic and universal testing procedure can not be defined, it can be expected that once hardware functionality is confirmed, attention shifts to software validation. Debugging software running on an autonomous system is in general more challenging than debugging common application software due to several specific factors, like the following ones.
\begin{itemize}
  \item It is more difficult to trace program execution on an autonomous, potentially mobile device: direct access to the system may not be available, and the tools at hand are often limited or not very sophisticated.
  \item In the event of issues, both the software and hardware in use must be checked as the problem may be due to hardware faults or incorrect use.
  \item Conducting tests is not as simple as running a program and verifying its results: the safety of the apparatus and its operators must be ensured throughout the entire procedure, which can make even the organization of a test a challenging task.
\end{itemize}

Given the difficulty and importance of this phase, it is reasonable to employ tools and methodologies that can facilitate these tasks as much as possible.
